\begin{tikzpicture}[
    fiber/.style={semgblue, line width=1.15pt},
    nerve/.style={semgorange, line width=1.1pt},
    label/.style={font=\small, align=center}
]
% Nơron vận động và axon
\draw[fill=semgorange!18, draw=semgorange, line width=1pt]
    (0,1.9) circle (0.48);
\node[label] at (0,1.9) {nơron\\vận động};
\draw[nerve] (0.45,1.9) -- (2.1,1.9);
\node[label, above] at (1.25,1.9) {axon};
% Các nhánh đến sợi cơ
\foreach \y in {0.3,0.95,1.6,2.25,2.9,3.55} {
    \draw[nerve] (2.1,1.9) -- (3.0,\y);
    \draw[fiber] (3.0,\y) -- (8.9,\y);
    \fill[semgorange] (3.0,\y) circle (1.7pt);
    \draw[-{Latex[length=1.5mm]}, semgred, line width=0.7pt]
        (5.1,\y+0.06) -- (5.85,\y+0.06);
}
\node[label, anchor=west] at (9.05,0.3) {sợi cơ 1};
\node[label, anchor=west] at (9.05,0.95) {sợi cơ 2};
\node[label, anchor=west] at (9.05,1.6) {sợi cơ 3};
\node[label, anchor=west] at (9.05,2.25) {sợi cơ 4};
\node[label, anchor=west] at (9.05,2.9) {sợi cơ 5};
\node[label, anchor=west] at (9.05,3.55) {sợi cơ 6};
% Chú thích vùng phát xung
\draw[dashed, semgred, line width=0.8pt] (3.0,0.02) -- (3.0,3.83);
\node[label, text=semgred, rotate=90] at (2.67,1.92) {khớp thần kinh--cơ};
\node[label, below=7mm of $(6.0,0.3)$, text width=8.5cm]
    {Một đơn vị vận động gồm một nơron vận động và toàn bộ các sợi cơ do nó chi phối.
    Mũi tên đỏ biểu diễn hướng lan truyền điện thế hoạt động trên màng sợi cơ.};
\end{tikzpicture}
